\chapter{Evaluación experimental}

\section{Configuración de los experimentos}
\subsection{Particiones de los datos}
\subsection{Hiperparámetros y métricas}

\section{Entrenamiento desde cero y transferencia desde MEG-XL}

\subsection{Inicialización desde cero y desde MEG-XL}

% La comparación entre inicialización preentrenada y desde cero es el eje experimental de la propuesta. La condición preentrenada carga parcialmente el checkpoint \texttt{meg-xl-med.ckpt}. Esta carga aprovecha los pesos compatibles del Transformer y de los componentes generales de la arquitectura, pero permite reinicializar las partes que no encajan con la configuración EEG. La condición desde cero utiliza la misma arquitectura y el mismo pipeline, pero sin cargar pesos previos.

% En la fase de lectura, la diferencia entre ambas condiciones es clara: una parte del checkpoint de MEG-XL y la otra comienza con pesos aleatorios. En la fase de escucha, ambas condiciones continúan desde su propio checkpoint de lectura. Por tanto, la comparación final no enfrenta simplemente MEG-XL contra un modelo aleatorio, sino dos trayectorias completas de adaptación:

% \begin{itemize}
% \setlength{\itemsep}{0pt}
% \setlength{\parsep}{0pt}
% \setlength{\topsep}{2pt}
% \setlength{\partopsep}{0pt}
% \item MEG-XL $\rightarrow$ lectura EEG $\rightarrow$ escucha EEG;
% \item inicialización aleatoria $\rightarrow$ lectura EEG $\rightarrow$ escucha EEG.
% \end{itemize}

% Este diseño controla el efecto del curriculum. Si ambos modelos mejoran por entrenar primero con lectura y después con escucha, la diferencia restante se debe a la inicialización. También permite detectar transferencia negativa: podría ocurrir que los pesos aprendidos en MEG introduzcan sesgos espaciales o estadísticos poco adecuados para EEG y ralenticen la adaptación frente al entrenamiento desde cero.


\section{Comparación de tokenizadores}

\section{Comparación de tareas}
\subsection{Lectura}
\subsection{Escucha}
\subsection{Lectura y escucha combinadas}

\section{Discusión de los resultados}
\subsection{Principales conclusiones experimentales}
\subsection{Limitaciones}

%%%%%%%%%%%%%%%%%%%%%%%%%%%%%%%%%%%%%%%%%%%%%%%%%%%%%%%%%%%%%%%%%%%%%%%%%%%%%%%
%                                 CONCLUSIONS                                 %
%%%%%%%%%%%%%%%%%%%%%%%%%%%%%%%%%%%%%%%%%%%%%%%%%%%%%%%%%%%%%%%%%%%%%%%%%%%%%%%

